% !TeX root = ../main.tex
% Original qualitative cutting example. Hatched caps are physical cut surfaces;
% the crossed dashed relation is an invalid old shared constraint, not a method
% result. Stable surface histories and newborn boundary histories differ.
\begingroup
% Shared local styles for the four original problem illustrations in Section 2.2.
% Loaded inside each figure's group; all meanings remain legible in grayscale.
\linespread{1}\selectfont
\tikzset{
  rpBase/.style={x=1cm,y=1cm,text=black,line cap=round,line join=round,
    font=\songti\fontsize{9.8pt}{12pt}\selectfont,
    every node/.style={align=center,inner sep=0pt,outer sep=0pt}},
  rpTitle/.style={font=\heiti\fontsize{10.3pt}{13pt}\selectfont},
  rpNote/.style={font=\songti\fontsize{9.3pt}{12pt}\selectfont},
  rpFrame/.style={draw=black,line width=0.65pt},
  rpArrow/.style={draw=black,line width=0.8pt,
    -{Latex[length=1.7mm,width=1.15mm]}},
  rpMotion/.style={draw=black,line width=0.95pt,
    -{Latex[length=1.6mm,width=1.1mm]}},
  rpAction/.style={draw=black,line width=1.5pt,
    -{Latex[length=2.2mm,width=1.7mm]}},
  rpHistory/.style={draw=black!65,line width=0.65pt,
    dash pattern=on 1.4pt off 1.6pt},
  rpFuture/.style={draw=black,line width=0.75pt,
    dash pattern=on 3pt off 2pt},
  rpBoundary/.style={draw=black,line width=1.25pt}
}
\newcommand{\rpPanel}[5]{%
  \draw[rpFrame] (#1,#2) rectangle (#3,#4);
  \path[fill=black!5] (#1+0.01,#4-0.55) rectangle (#3-0.01,#4-0.01);
  \draw[rpFrame] (#1,#4-0.55) -- (#3,#4-0.55);
  \node[rpTitle] at ({(#1+#3)/2},{#4-0.275}) {#5};
}
\newcommand{\rpFooter}[1]{%
  \draw[rpFrame] (0.10,0.10) rectangle (15.90,0.77);
  \node at (8,0.435) {#1};
}
% Hatching is explicitly drawn and clipped, so no additional TikZ libraries
% or color meanings are required. Argument: a closed local-coordinate path.
\newcommand{\rpHatchFace}[1]{%
  \path[fill=black!5] #1;
  \begin{scope}
    \clip #1;
    \foreach \rhx in {-6,-5.8,...,8}{%
      \draw[black,line width=0.35pt] (\rhx,-1) -- ++(3,6);
    }
  \end{scope}
  \draw[rpBoundary] #1;
}

\newcommand{\rpSplitBody}{%
  % Two pieces face the opening so both newborn cut surfaces can be seen.
  \path[draw=black,fill=white,line width=0.65pt]
    (0,1.10) -- (0.40,1.45) -- (1.50,1.45) -- (1.10,1.10) -- cycle;
  \path[draw=black,fill=black!5,line width=0.65pt] (0,0) rectangle (1.10,1.10);
  \rpHatchFace{(1.10,0) -- (1.50,0.35) -- (1.50,1.45) -- (1.10,1.10) -- cycle}
  \path[draw=black,fill=white,line width=0.65pt]
    (3.00,1.10) -- (2.60,1.45) -- (3.70,1.45) -- (4.10,1.10) -- cycle;
  \path[draw=black,fill=black!5,line width=0.65pt] (3.00,0) rectangle (4.10,1.10);
  \rpHatchFace{(3.00,0) -- (2.60,0.35) -- (2.60,1.45) -- (3.00,1.10) -- cycle}
  \fill[black] (0.55,0.55) circle[radius=0.055];
  \fill[black] (3.55,0.55) circle[radius=0.055];
}
\begin{tikzpicture}[rpBase]
  \path[use as bounding box] (0,0) rectangle (16,5.35);
  \rpPanel{0.10}{1.05}{5.00}{5.20}{结构事件前}
  \rpPanel{5.55}{1.05}{10.45}{5.20}{结构事件后}
  \rpPanel{11.00}{1.05}{15.90}{5.20}{预期表征能力}

  \begin{scope}[shift={(0.90,2.35)}]
    \path[draw=black,fill=white,line width=0.65pt]
      (0,1.10) -- (0.40,1.45) -- (3.10,1.45) -- (2.70,1.10) -- cycle;
    \path[draw=black,fill=black!5,line width=0.65pt] (0,0) rectangle (2.70,1.10);
    \path[draw=black,fill=black!10,line width=0.65pt]
      (2.70,0) -- (3.10,0.35) -- (3.10,1.45) -- (2.70,1.10) -- cycle;
    \draw[rpHistory] (1.35,0) -- (1.35,1.10) -- (1.75,1.45);
    \fill[black] (0.55,0.55) circle[radius=0.055];
    \fill[black] (2.15,0.55) circle[radius=0.055];
  \end{scope}
  \draw[rpFrame] (1.35,4.03) rectangle (3.75,4.55);
  \node[rpNote] at (2.55,4.29) {局部切割};
  \draw[draw=black,line width=0.55pt] (2.55,4.03) -- (2.25,3.35);
  \draw[rpFrame] (0.40,1.30) rectangle (4.70,1.85);
  \node[rpNote] at (2.55,1.575) {原有表面和连接};
  \draw[rpArrow] (5.07,3.05) -- (5.48,3.05);

  \begin{scope}[shift={(5.95,2.35)}]
    \rpSplitBody
    \draw[rpHistory,line width=0.85pt] (1.52,0.65) -- (2.58,0.65);
    \draw[black,line width=1.0pt]
      (1.90,0.50) -- (2.20,0.80) (1.90,0.80) -- (2.20,0.50);
  \end{scope}
  \draw[rpFrame] (6.80,4.03) rectangle (9.20,4.55);
  \node[rpNote] at (8.00,4.29) {物理新生表面};
  \draw[draw=black,line width=0.55pt]
    (7.65,4.03) -- (7.25,3.30) (8.35,4.03) -- (8.75,3.30);
  \draw[rpFrame] (5.85,1.30) rectangle (10.15,1.85);
  \node[rpNote] at (8.00,1.575) {旧共享约束可能失效};
  \draw[rpArrow] (10.52,3.05) -- (10.93,3.05);

  \begin{scope}[shift={(11.40,2.35)}]
    \rpSplitBody
    \draw[rpHistory] (0.15,0.25) .. controls (0.20,0.45) and (0.35,0.55) .. (0.55,0.55);
    \draw[rpHistory] (3.15,0.25) .. controls (3.20,0.45) and (3.35,0.55) .. (3.55,0.55);
  \end{scope}
  \draw[rpFrame] (11.30,4.03) rectangle (15.60,4.55);
  \node[rpNote] at (13.45,4.29) {稳定区域延续有效历史};
  \draw[rpFrame] (11.30,1.30) rectangle (15.60,1.85);
  \node[rpNote] at (13.45,1.575) {新生边界从形成时刻记录};

  \rpFooter{研究关注：如何表达局部结构改变并延续有观测支持的历史？}
\end{tikzpicture}
\endgroup
